\documentclass[12pt, a4paper]{article}

% ── PACKAGES ──────────────────────────────────────────────────
\usepackage[utf8]{inputenc}
\usepackage[T1]{fontenc}
\usepackage{geometry}
\geometry{margin=2.5cm}
\usepackage{hyperref}
\hypersetup{
    colorlinks=true,
    linkcolor=blue,
    urlcolor=blue,
    citecolor=blue,
    pdftitle={Anti-TIGIT Before Anti-PD-1 in
              Claudin-Low Memory-Low Breast Cancer:
              Patient Selection by Lineage Memory
              Score and FOXP3/CD8A Ratio as the
              Solution to Anti-TIGIT Trial Failures},
    pdfauthor={Eric Robert Lawson},
    pdfkeywords={anti-TIGIT, tiragolumab, TIGIT,
                 claudin-low, memory-low, FOXP3, CD8A,
                 Treg, checkpoint inhibitor, anti-PD-1,
                 atezolizumab, SKYLINE, NCT06175390,
                 belrestotug, breast cancer, triple-negative,
                 patient selection, FOXA1, SPDEF, GATA3,
                 CT antigen, attractor geometry,
                 Waddington landscape, OrganismCore,
                 GSE176078, Type 4, root lock,
                 immune compartment, precision oncology}
}
\usepackage{amsmath}
\usepackage{booktabs}
\usepackage{longtable}
\usepackage{array}
\usepackage{parskip}
\usepackage{microtype}
\usepackage{authblk}
\usepackage{titlesec}
\usepackage{fancyhdr}
\usepackage{xcolor}
\usepackage{float}
\usepackage{enumitem}
\usepackage{setspace}
\usepackage{multirow}

% ── COLOURS ───────────────────────────────────────────────────
\definecolor{repobox}{RGB}{240,247,255}
\definecolor{findingbox}{RGB}{245,255,245}
\definecolor{convergentbox}{RGB}{255,250,235}
\definecolor{mechanismbox}{RGB}{250,245,255}
\definecolor{novelbox}{RGB}{255,248,220}
\definecolor{actionbox}{RGB}{230,255,230}
\definecolor{urgentbox}{RGB}{255,235,235}
\definecolor{failurebox}{RGB}{255,240,240}
\definecolor{statbox}{RGB}{240,240,255}

% ── HEADER / FOOTER ───────────────────────────────────────────
\pagestyle{fancy}
\fancyhf{}
\lhead{\footnotesize CS-LIT-19/20: Anti-TIGIT Sequence
       in Claudin-Low / Memory-Low}
\rhead{\footnotesize OrganismCore \textbar{} 2026-03-06}
\rfoot{\small Page \thepage}
\renewcommand{\headrulewidth}{0.4pt}
\setlength{\headheight}{24pt}
\addtolength{\topmargin}{-10pt}

% ── SECTION FORMATTING ────────────────────────────────────────
\titleformat{\section}
  {\large\bfseries}{\thesection}{1em}{}[\titlerule]
\titleformat{\subsection}
  {\normalsize\bfseries}{\thesubsection}{1em}{}
\titleformat{\subsubsection}
  {\small\bfseries}{\thesubsubsection}{1em}{}

% ══════════════════════════════════════════════════════════════
% TITLE
% ══════════════════════════════════════════════════════════════
\title{
    \vspace{-1em}
    \textbf{Anti-TIGIT Before Anti-PD-1 in Claudin-Low,
    Memory-Low Breast Cancer:}\\[0.5em]
    \large Patient Selection by Lineage Memory Score
    and FOXP3/CD8A Ratio as the Mechanistic Solution
    to Anti-TIGIT Trial Failures\\[0.4em]
    \normalsize \textit{OrganismCore Framework ---
    Documents CS-LIT-19 and CS-LIT-20}\\[0.3em]
    \small Derived from Waddington Attractor Geometry
    of 19,542 Single Cancer Cells (GSE176078)\\[0.3em]
    \small Confirmed by Taylor et al.\ \textit{JCI} 2017
    and Pommier et al.\ \textit{Nature Communications} 2020
}

\author[1]{Eric Robert Lawson}
\affil[1]{
    OrganismCore\\
    \href{mailto:OrganismCore@proton.me}
         {OrganismCore@proton.me}\\
    ORCID:
    \href{https://orcid.org/0009-0002-0414-6544}
         {0009-0002-0414-6544}
}

\date{March 6, 2026}

% ══════════════════════════════════════════════════════════════
\begin{document}
% ══════════════════════════════════════════════════════════════

\maketitle
\thispagestyle{fancy}

% ── REPOSITORY POINTER ────────────────────────────────────────
\begin{center}
\colorbox{repobox}{
\begin{minipage}{0.88\textwidth}
\vspace{0.5em}
\centering
\textbf{Full computational record --- scripts, data
caches, reasoning artifacts, and raw logs:}\\[0.4em]
\href{https://github.com/Eric-Robert-Lawson/attractor-oncology}
{\texttt{https://github.com/Eric-Robert-Lawson/attractor-oncology}}
\\[0.3em]
\small Repository ID: 1172048282 \textbar{}
MIT License \textbar{}
Python \textbar{}
Public\\[0.2em]
Part of the OrganismCore BRCA Cross-Subtype Analysis
(BRCA-S8h, 2026-03-05).
All predictions locked before literature review.
\vspace{0.5em}
\end{minipage}}
\end{center}

\vspace{0.8em}

% ── TIME-SENSITIVE BOX ──��─────────────────────────────────────
\begin{center}
\colorbox{urgentbox}{
\begin{minipage}{0.88\textwidth}
\vspace{0.5em}
\centering
\textbf{TIME-SENSITIVE NOTE}\\[0.4em]
\small
SKYLINE (NCT06175390) --- Phase~II tiragolumab
$+$ atezolizumab $+$ chemotherapy in TNBC, Institut
Curie. PI: Prof.\ Fran\c{c}ois-Cl\'{e}ment Bidard.
Enrollment open; estimated completion: July 2026.\\[0.3em]
The trial includes a \textbf{multi-omics biomarker arm}.
It does \textbf{not} currently use claudin-low subtype,
memory-low lineage score, or FOXP3/CD8A ratio as
patient stratification or correlative variables.\\[0.3em]
These variables could be added to the biomarker
analysis on already-planned tissue collection
without protocol amendment.\\[0.4em]
\textbf{Belrestotug (GSK/iTeos) discontinued May 2025}
after Phase~II failures in PD-L1-high NSCLC and HNSCC.
No claudin-low enrichment. No memory-low selection.
No FOXP3/CD8A stratification. The failures are
\textit{consistent with the framework}: unselected
populations are not the predicted beneficiary.
\vspace{0.5em}
\end{minipage}}
\end{center}

\vspace{0.8em}

% ══════════════════════════════════════════════════════════════
\begin{abstract}
% ══════════════════════════════════════════════════════════════
\noindent
Anti-TIGIT immunotherapy has failed in multiple
Phase~II/III trials in unselected cancer populations
(most recently belrestotug, May 2025, in PD-L1-high
NSCLC and HNSCC). The OrganismCore Waddington
attractor geometry framework, applied to 19,542 single
cancer cells (GSE176078), identifies the specific
patient population that is predicted to benefit from
anti-TIGIT and the reason why unselected trials fail:
the predicted beneficiary is the claudin-low,
memory-low breast cancer subgroup specifically, and
the sequence is required --- anti-TIGIT must precede
anti-PD-1, not accompany it.

\textbf{Finding CS-LIT-19 (the sequence):}
Claudin-low (CL) breast cancer is the Type~4 root
lock attractor --- no luminal identity programme,
highest CT antigen load, highest immune infiltration,
most stem-like. Within CL, the memory-low subgroup
(FOXA1/SPDEF/GATA3 absent; Pommier et al.\
\textit{Nature Communications} 2020) is the most
deeply arrested and most immune-infiltrated
subpopulation. Taylor et al.\ (\textit{JCI} 2017,
Morel group) demonstrated that anti-PD-1 alone in
claudin-low breast cancer amplifies Tregs before
depleting them, making outcomes worse when given
first. Only rigorous Treg depletion produces tumour
growth delay. Anti-TIGIT (tiragolumab, ociperlimab)
depletes Tregs via the TIGIT/CD155 co-inhibitory
axis. The required sequence is therefore: anti-TIGIT
first (Treg depletion) $\rightarrow$ anti-PD-1 second
(checkpoint release with Tregs cleared). Reversing
the sequence or using anti-PD-1 without prior Treg
depletion worsens outcomes in CL.

\textbf{Finding CS-LIT-20 (the patient selection
instrument):}
The FOXP3/CD8A ratio is the strongest single
immune biomarker in claudin-low breast cancer
(HR$=2.212$ in CL-specific survival analysis from
GSE176078). High FOXP3/CD8A means Tregs dominate
cytotoxic T cells --- the exact condition that
anti-TIGIT addresses. The memory-low lineage score
(FOXA1/SPDEF/GATA3 absent) identifies the
subpopulation within CL with the highest CT
antigen load and highest immune infiltration ---
the patients for whom the immune compartment is
the primary and potentially actionable target.

\textbf{Why unselected trials fail:}
Belrestotug was tested in PD-L1-high NSCLC and
HNSCC --- neither is claudin-low breast cancer.
Neither trial used FOXP3/CD8A stratification.
Neither specified the anti-TIGIT first sequence as
a protocol requirement. Anti-PD-1 in unselected
populations amplifies Tregs as a resistance mechanism.
The failures are mechanistically predictable from
the framework and do not undermine the prediction
for the specific selected population.

NCT06175390 (SKYLINE; tiragolumab $+$ atezolizumab
in TNBC; Institut Curie; PI: Prof.\ Bidard;
enrolling until July 2026) has a multi-omics
biomarker arm that provides an immediate opportunity
to add claudin-low subtype and FOXP3/CD8A ratio as
stratification variables without protocol amendment.

Verdict: CONVERGENT-NOVEL.
\end{abstract}

\vspace{0.5em}
\noindent\textbf{Keywords:}
anti-TIGIT; tiragolumab; TIGIT; claudin-low;
memory-low; FOXP3; CD8A; Treg; checkpoint inhibitor;
anti-PD-1; atezolizumab; SKYLINE; NCT06175390;
belrestotug; breast cancer; patient selection;
FOXA1; SPDEF; GATA3; CT antigen; attractor geometry;
Waddington landscape; OrganismCore; Type~4; root lock;
immune compartment; precision oncology;
CONVERGENT-NOVEL

\newpage
\tableofcontents
\newpage

% ══════════════════════════════════════════════════════════════
\section{Document Metadata}
% ══════════════════════════════════════════════════════════════

\begin{center}
\begin{tabular}{ll}
\toprule
\textbf{Field} & \textbf{Value} \\
\midrule
Document IDs & CS-LIT-19 and CS-LIT-20 \\
Parent document & BRCA-S8h (2026-03-05) \\
Verdict & CONVERGENT-NOVEL \\
Date & 2026-03-06 \\
Author & Eric Robert Lawson \\
ORCID & 0009-0002-0414-6544 \\
Organisation & OrganismCore \\
Contact & OrganismCore@proton.me \\
Repository &
\href{https://github.com/Eric-Robert-Lawson/attractor-oncology}
{attractor-oncology} (ID: 1172048282) \\
Derivation data & GSE176078 (19,542 single cancer cells) \\
Key supporting references &
Taylor et al.\ \textit{JCI} 2017 \\
& Pommier et al.\ \textit{Nat Comms} 2020 \\
Target trial &
NCT06175390 (SKYLINE; Institut Curie) \\
Trial PI &
Prof.\ Fran\c{c}ois-Cl\'{e}ment Bidard \\
CS-LIT-2 DOI &
\href{https://doi.org/10.5281/zenodo.18884158}
{10.5281/zenodo.18884158} \\
CS-LIT-6 DOI &
\href{https://doi.org/10.5281/zenodo.18891770}
{10.5281/zenodo.18891770} \\
License & CC BY 4.0 \\
Biological contradictions & 0 \\
\bottomrule
\end{tabular}
\end{center}

% ══════════════════════════════════════════════════════════════
\section{Claudin-Low: The Type~4 Root Lock Attractor}
% ══════════════════════════════════════════════════════════════

\subsection{Position in the six lock type classification}

In the OrganismCore six lock type classification
(CS-LIT-2,
\href{https://doi.org/10.5281/zenodo.18884158}
{DOI: 10.5281/zenodo.18884158}), claudin-low (CL)
breast cancer is the \textbf{Type~4 root lock}:

\begin{center}
\begin{tabular}{p{3cm}p{9cm}}
\toprule
\textbf{Feature} & \textbf{Description} \\
\midrule
Attractor type &
Type~4: Root lock. Cell arrested at the progenitor
root. The luminal identity programme was never
activated --- it is absent, not silenced. \\[0.4em]
FOXA1/EZH2 ratio &
$0.10$ --- lowest of all measurable subtypes.
FOXA1 is near-zero (no luminal programme).
EZH2 is elevated (the cell never committed). \\[0.4em]
What this means &
No FOXA1, GATA3, ESR1, or luminal circuit to
target with endocrine therapy. The identity
compartment cannot be the primary target.
\textbf{The immune compartment is the only
accessible therapeutic axis.} \\[0.4em]
CL depth score &
EZH2-free PCA: $6.572$ --- deepest subtype,
exceeding TNBC ($6.063$) when measured
without EZH2 distortion. \\
\bottomrule
\end{tabular}
\end{center}

\subsection{The memory-low subgroup (Pommier et al.\
\textit{Nature Communications} 2020)}

Within claudin-low, Pommier et al.\ (2020) identified
a memory-low subgroup (subgroup~1) defined by:

\begin{itemize}
    \item \textbf{FOXA1/SPDEF/GATA3 absent:}
    No luminal lineage memory markers. The cell has
    not retained any transcriptional memory of
    luminal commitment.
    \item \textbf{Highest CT antigen load:}
    Cancer-testis antigens (MAGE-A1, NY-ESO-1,
    SSX family) are de-repressed in deeply
    undifferentiated cells --- epigenetic
    de-repression of germline programmes. This
    is the direct consequence of the root lock
    geometry: cells at the progenitor root have
    not silenced germline programmes.
    \item \textbf{Highest immune infiltration:}
    CT antigens are immunogenic. High CT antigen
    expression drives immune cell recruitment.
    The memory-low subgroup has the highest
    immune infiltration of any CL subgroup.
    \item \textbf{Most stem-like:}
    The most undifferentiated subgroup within
    the most undifferentiated subtype. The deepest
    point in the framework's attractor landscape.
\end{itemize}

The memory-low subgroup is the primary predicted
beneficiary of immune therapy --- not CL as a
whole, and certainly not unselected TNBC.

% ══════════════════════════════════════════════════════════════
\section{Finding CS-LIT-19: The Required Sequence}
% ══════════════════════════════════════════════════════════════

\subsection{Why anti-PD-1 alone fails in CL}

Taylor et al.\ (\textit{Journal of Clinical
Investigation}, 2017, Morel group) is the
definitive experimental paper on immune therapy
in claudin-low breast cancer. The key findings:

\begin{center}
\colorbox{convergentbox}{
\begin{minipage}{0.88\textwidth}
\vspace{0.5em}
\textbf{TAYLOR et al.\ JCI 2017 --- KEY FINDINGS:}\\[0.4em]
\begin{enumerate}[noitemsep]
    \item Anti-PD-1 alone in claudin-low mouse
    mammary tumours \textbf{amplifies Tregs}
    before depleting them --- the initial
    response to PD-1 blockade in CL is Treg
    expansion, not Treg depletion.
    \item Giving anti-PD-1 first therefore
    \textbf{worsens the immune suppression
    environment} before potentially improving it.
    Outcomes are worse when anti-PD-1 precedes
    Treg depletion.
    \item \textbf{Rigorous Treg depletion first}
    (depleting FOXP3$^+$ cells) produces
    substantial tumour growth delay in CL
    models.
    \item The mechanism: in CL, Tregs dominate
    the immune environment via FOXP3 expression
    and TIGIT/CD155 co-inhibitory signalling.
    CD8$^+$ cytotoxic T cells are present but
    are suppressed by Treg activity. Anti-PD-1
    releases PD-1-suppressed T cells globally,
    including Tregs --- enhancing Treg activity
    before effector T cell activity.
\end{enumerate}
\vspace{0.5em}
\end{minipage}}
\end{center}

\subsection{Why anti-TIGIT specifically}

Anti-TIGIT (tiragolumab, ociperlimab, vibostolimab)
targets the TIGIT/CD155 co-inhibitory axis. TIGIT
is expressed on Tregs and on exhausted CD8$^+$ T cells.
CD155 (PVR) is expressed on tumour cells and antigen-
presenting cells. The TIGIT/CD155 interaction:

\begin{itemize}
    \item Suppresses Treg-killing of tumour cells
    (Tregs use TIGIT to maintain dominant
    suppression)
    \item Maintains CD8$^+$ T cell exhaustion
    \item Anti-TIGIT disrupts both: it partially
    depletes Tregs (via Fc-mediated ADCC of
    TIGIT-expressing Tregs) and releases CD8$^+$
    T cells from TIGIT-mediated exhaustion
\end{itemize}

Anti-TIGIT is the mechanistically correct Treg-depleting
agent for the CL context because TIGIT is the
co-inhibitory molecule through which Tregs dominate
the CL immune microenvironment.

\subsection{The required sequence}

\begin{center}
\colorbox{mechanismbox}{
\begin{minipage}{0.88\textwidth}
\vspace{0.5em}
\textbf{THE REQUIRED SEQUENCE IN CL MEMORY-LOW:}\\[0.5em]
\textbf{Step~1:} Anti-TIGIT (tiragolumab or
equivalent)\\
$\downarrow$\\
TIGIT/CD155 axis disrupted\\
Tregs depleted (ADCC on TIGIT$^+$ Tregs)\\
CD8$^+$ T cell TIGIT-exhaustion partially reversed\\
Immune microenvironment shifted from
Treg-dominant to effector-permissive\\[0.5em]
\textbf{Step~2:} Anti-PD-1 or anti-PD-L1
(pembrolizumab, atezolizumab)\\
$\downarrow$\\
PD-1/PD-L1 axis disrupted in a Treg-cleared
environment\\
CD8$^+$ cytotoxic T cells activated against
CT antigen-expressing tumour cells\\
Durable anti-tumour immune response\\[0.5em]
\textbf{If reversed (anti-PD-1 first):}\\
Anti-PD-1 releases PD-1 brake globally\\
Tregs expand (Taylor JCI 2017)\\
Immune suppression worsens before improving\\
The window for anti-TIGIT Treg depletion is
compromised by prior Treg expansion
\vspace{0.5em}
\end{minipage}}
\end{center}

\noindent The sequence is a protocol requirement, not
a preference. The biology does not support
administering these agents simultaneously or in
reverse order in the CL memory-low population.

% ══════════════════════════════════════════════════════════════
\section{Finding CS-LIT-20: The Patient Selection
Instrument --- FOXP3/CD8A Ratio}
% ══════════════════════════════════════════════════════════════

\subsection{The ratio and its biological meaning}

\begin{center}
\colorbox{statbox}{
\begin{minipage}{0.88\textwidth}
\vspace{0.5em}
\centering
\textbf{FOXP3/CD8A RATIO IN CLAUDIN-LOW}\\[0.5em]
\textbf{HR = 2.212 in CL-specific survival analysis
(GSE176078)}\\[0.4em]
High FOXP3/CD8A $\rightarrow$ Tregs dominate
cytotoxic T cells\\
Low FOXP3/CD8A $\rightarrow$ Cytotoxic T cells
relatively dominant\\[0.3em]
In CL, high FOXP3/CD8A patients carry HR$=2.212$
for worse survival --- confirming that Treg
dominance is the primary driver of poor outcomes
in this subtype.\\[0.3em]
\textbf{Anti-TIGIT directly addresses high
FOXP3/CD8A: it depletes FOXP3$^+$ Tregs and
releases CD8$^+$ T cells.}
\vspace{0.5em}
\end{minipage}}
\end{center}

\subsection{How to use the ratio clinically}

The FOXP3/CD8A ratio is computable from:
\begin{itemize}
    \item \textbf{IHC:} FOXP3 IHC (standard clinical
    antibody, available in any oncology pathology lab)
    $+$ CD8A IHC. H-score ratio or cell-count ratio.
    \item \textbf{RNA-seq:} FOXP3 transcript /
    CD8A transcript from bulk or single-cell data.
    \item \textbf{NanoString / immune panel:}
    Computable from any multi-gene immune panel
    that includes both markers.
\end{itemize}

The prediction: within CL (identified by lineage
memory score: FOXA1/SPDEF/GATA3 absent), patients
with the highest FOXP3/CD8A ratio derive the greatest
benefit from anti-TIGIT first $\rightarrow$ anti-PD-1
second. The ratio is a continuous predictor of
benefit magnitude, not a binary cut-point.

\subsection{The two-variable patient selection
instrument}

Full patient selection for the predicted beneficiary
population requires both variables:

\begin{center}
\begin{tabular}{p{4cm}p{7.5cm}}
\toprule
\textbf{Variable} & \textbf{Identification} \\
\midrule
\textbf{CL subtype + memory-low} &
Lineage memory score: FOXA1, SPDEF, GATA3 all absent
(IHC H-score $<10$ for each). This identifies the
Type~4 root lock attractor within TNBC.
Alternatively: PAM50 claudin-low call $+$
low FOXA1/EZH2 ratio (CS-LIT-22,
\href{https://doi.org/10.5281/zenodo.18892788}
{DOI: 10.5281/zenodo.18892788}). \\[0.5em]
\textbf{FOXP3/CD8A ratio} &
IHC or RNA-based ratio. Quantifies Treg dominance.
The patient selection confirms the Treg-dominant
immune environment that anti-TIGIT addresses.
High ratio ($>$ median in CL cohort) identifies
the predicted primary beneficiary. \\
\bottomrule
\end{tabular}
\end{center}

% ══════════════════════════════════════════════════════════════
\section{The Anti-TIGIT Trial Failures: Why They
Are Consistent with This Framework}
% ══════════════════════════════════════════════════════════════

\begin{center}
\colorbox{failurebox}{
\begin{minipage}{0.88\textwidth}
\vspace{0.5em}
\textbf{BELRESTOTUG (GSK/iTeos): DISCONTINUED
MAY 2025}\\[0.4em]
\textbf{Trials:}\\
GALAXIES Lung-201 (NCT05565378): belrestotug
$+$ dostarlimab (anti-PD-1) vs dostarlimab alone
in PD-L1-high ($\geq50\%$ TPS) unresectable
locally advanced or metastatic NSCLC.\\[0.3em]
GALAXIES H\&N-202 (NCT06062420): belrestotug
$+$ dostarlimab in PD-L1-positive recurrent or
metastatic HNSCC.\\[0.4em]
\textbf{Result:} Both trials failed to show
meaningful improvement in PFS over dostarlimab
monotherapy. ORR improved in NSCLC but not to
clinical significance. Program discontinued.\\[0.4em]
\textbf{Framework interpretation:}\\
\begin{enumerate}[noitemsep]
    \item Neither trial enrolled claudin-low
    breast cancer. NSCLC and HNSCC are not
    predicted beneficiary populations.
    \item Neither trial used FOXP3/CD8A
    stratification. Treg-dominant patients
    were not selected.
    \item Neither trial specified anti-TIGIT
    first as a sequence requirement ---
    both gave concurrent combination.
    \item PD-L1-high selection is not the
    same as CL memory-low selection: PD-L1
    measures checkpoint expression, not the
    Treg/CD8 dominance ratio that determines
    anti-TIGIT benefit.
\end{enumerate}
\textbf{The failures are exactly what the
framework predicts for unselected, non-CL,
wrong-sequence populations.}\\[0.3em]
They do not contradict the prediction for
memory-low CL with high FOXP3/CD8A given
anti-TIGIT first.
\vspace{0.5em}
\end{minipage}}
\end{center}

\begin{center}
\begin{tabular}{p{4.5cm}p{3.5cm}p{3.5cm}}
\toprule
\textbf{Trial characteristic} &
\textbf{Belrestotug} &
\textbf{Framework prediction} \\
\midrule
Tumour type &
NSCLC, HNSCC &
CL breast cancer \\[0.3em]
Patient selection &
PD-L1 $\geq50\%$ TPS &
Memory-low $+$ FOXP3/CD8A high \\[0.3em]
Sequence &
Concurrent combination &
Anti-TIGIT first \\[0.3em]
Treg stratification &
None &
Required \\[0.3em]
Result &
Failure &
Predicted failure (wrong population,
wrong sequence) \\
\bottomrule
\end{tabular}
\end{center}

% ══════════════════════════════════════════════════════════════
\section{SKYLINE (NCT06175390): The Immediate
Opportunity}
% ═════════════════════��════════════════════════════════════════

\begin{center}
\colorbox{actionbox}{
\begin{minipage}{0.88\textwidth}
\vspace{0.5em}
\textbf{SKYLINE TRIAL DETAILS:}\\[0.4em]
\textbf{Full title:} Phase~II study evaluating
tiragolumab $+$ atezolizumab $+$ chemotherapy
in TNBC.\\[0.3em]
\textbf{Sponsor:} Institut Curie, Paris.\\[0.3em]
\textbf{PI:} Prof.\ Fran\c{c}ois-Cl\'{e}ment Bidard.\\[0.3em]
\textbf{Contact:} Institut Curie, 26 rue d'Ulm,
75005 Paris, France. Tel: $+$33\,1\,44\,32\,46\,70.\\[0.3em]
\textbf{Cohorts:} Neoadjuvant (early TNBC) and
metastatic TNBC.\\[0.3em]
\textbf{Enrollment:} Ongoing. Estimated
completion: July 2026.\\[0.3em]
\textbf{Biomarker arm:} Multi-omics biomarker
collection included in trial design.\\[0.4em]
\textbf{Current gap:}\\
No claudin-low subtype annotation.\\
No memory-low (FOXA1/SPDEF/GATA3) stratification.\\
No FOXP3/CD8A ratio analysis.\\[0.4em]
\textbf{Proposed addition (no protocol amendment
needed):}\\
CL subtype annotation from tumour RNA or IHC
on already-planned biopsy tissue. FOXP3 and
CD8A IHC on the same tissue. This adds three
IHC stains (FOXA1, FOXP3, CD8A) and/or an RNA
panel component. No additional patient visit.\\[0.4em]
\textbf{The testable prediction in SKYLINE:}\\
Within the TNBC cohort, patients classified as
claudin-low memory-low with high FOXP3/CD8A
ratio show substantially better response to
tiragolumab $+$ atezolizumab than the
non-CL TNBC patients. If the sequence is
concurrently administered (as in the current
protocol) rather than anti-TIGIT first, the
prediction is attenuated but the direction
should still hold in this enriched subgroup.
\vspace{0.5em}
\end{minipage}}
\end{center}

% ══════════════════════════════════════════════════════════════
\section{What Is and Is Not in the Published
Literature}
% ══════════════════════════════════════════════════════════════

\subsection{What IS established}

\begin{enumerate}
    \item \textbf{Taylor et al.\ JCI 2017:}
    Anti-PD-1 amplifies Tregs first in CL.
    Treg depletion required before checkpoint
    release in CL. This is the exact mechanistic
    basis for the anti-TIGIT first sequence.

    \item \textbf{Pommier et al.\ Nat Comms 2020:}
    Memory-low CL subgroup (FOXA1/SPDEF/GATA3
    absent) characterised. Highest CT antigen
    load. Highest immune infiltration. Most
    stem-like.

    \item \textbf{FOXP3/CD8A imbalance in CL:}
    Taylor 2017 and Pommier 2020 both describe
    FOXP3-dominant, CD8-suppressed immune
    microenvironment in CL as the mechanistic
    basis for immune therapy resistance.

    \item \textbf{TIGIT biology:}
    TIGIT expressed on Tregs and exhausted
    CD8$^+$ T cells. Anti-TIGIT Fc-mediated
    ADCC of TIGIT$^+$ Tregs established.
    TIGIT/CD155 axis in immune suppression:
    published across multiple tumour types.

    \item \textbf{CL Type~4 root lock:}
    CS-LIT-2 (DOI: 10.5281/zenodo.18884158).
    EZH2-free PCA depth 6.572 (deepest subtype).
    No luminal identity programme to target.
    Immune compartment is the primary axis.
\end{enumerate}

\subsection{What is NOT in the literature}

\begin{center}
\colorbox{novelbox}{
\begin{minipage}{0.88\textwidth}
\vspace{0.5em}
\textbf{The following have no published equivalent
as of 2026-03-06:}
\begin{itemize}[noitemsep]
    \item Any clinical trial selecting CL patients
    for anti-TIGIT by claudin-low subtype.
    \item Any trial using memory-low (FOXA1/SPDEF/GATA3
    absent) as an enrichment criterion for
    anti-TIGIT eligibility.
    \item Any trial specifying anti-TIGIT first
    $\rightarrow$ anti-PD-1 second as a protocol
    requirement (rather than concurrent).
    \item Any trial or paper using FOXP3/CD8A
    ratio as a patient selection variable for
    anti-TIGIT with HR$=2.212$ in CL-specific
    survival analysis.
    \item The mechanistic explanation linking
    the belrestotug failures (wrong population,
    wrong sequence) to the CL patient selection
    framework.
    \item The two-variable patient selection
    instrument: memory-low lineage score $+$
    FOXP3/CD8A ratio.
\end{itemize}
\vspace{0.5em}
\end{minipage}}
\end{center}

% ══════════════════════════════════════════════════════════════
\section{Locked Statement}
% ═════════════════════════════════════════════════════════���════

\begin{center}
\colorbox{findingbox}{
\begin{minipage}{0.88\textwidth}
\vspace{0.6em}
\textbf{Stated once, precisely, without inflation:}
\\[0.5em]
Anti-TIGIT immunotherapy has failed in unselected
cancer populations because the predicted beneficiary
is specifically the claudin-low memory-low subgroup
(FOXA1/SPDEF/GATA3 absent; highest CT antigen load;
Pommier et al.\ Nat Comms 2020), the sequence
matters (anti-TIGIT must precede anti-PD-1; Taylor
et al.\ JCI 2017), and the patient selection
instrument is the FOXP3/CD8A ratio (HR$=2.212$ in
CL-specific survival analysis, GSE176078) not PD-L1
TPS. Belrestotug failed in PD-L1-high NSCLC and
HNSCC (May 2025) --- non-CL tumour types, no Treg
stratification, concurrent rather than sequential
dosing. These failures are mechanistically predicted
by the framework and do not undermine the prediction
for memory-low CL with high FOXP3/CD8A given
anti-TIGIT first. SKYLINE (NCT06175390; tiragolumab
$+$ atezolizumab $+$ chemotherapy in TNBC; Institut
Curie; PI: Prof.\ Bidard; enrollment open until
July 2026) has a multi-omics biomarker arm in which
CL subtype annotation and FOXP3/CD8A ratio could be
added without protocol amendment. No clinical trial
has selected patients by CL subtype for anti-TIGIT,
specified the anti-TIGIT first sequence as a protocol
requirement, or used FOXP3/CD8A as a selection
variable. Verdict: CONVERGENT-NOVEL.
\vspace{0.5em}
\end{minipage}}
\end{center}

\end{document}
